\section{Interfaz de usuario del panel de control}
\label{anexo:interfaz}

El sistema se opera a través de un panel de control web construido con Next.js que
actúa como capa de presentación sobre la API REST (Sección~\ref{sec:infraestructura}).
Su diseño persigue un objetivo concreto, alineado con el relato de este trabajo:
mostrar de un vistazo \emph{dónde razona un agente, dónde ejecuta código determinista
y dónde se solicita aprobación humana}. Por ello la interfaz no es una consola de
depuración, sino un panel de control de MLOps con un lenguaje visual deliberadamente
sobrio --- fondo neutro (\textit{zinc}), un único color de acento (\textit{indigo}) y
tipografía monoespaciada para identificadores y métricas, de modo que estos resalten
sin recurrir a adornos.

La aplicación se organiza en tres ventanas accesibles desde la barra de navegación
superior, cada una asociada a una fase distinta del ciclo de vida del modelo:

\begin{itemize}
  \item \textbf{Pipeline} --- lanzamiento y seguimiento en vivo de una ejecución,
        incluidas las puertas de aprobación humana.
  \item \textbf{Experiments} --- consulta de las ejecuciones de entrenamiento ya
        registradas y sus métricas, leídas de MLflow.
  \item \textbf{Observability} --- observabilidad agregada del sistema: salud del
        \textit{pipeline}, actividad de los LLM y su coste estimado, y uso de
        herramientas.
\end{itemize}

A continuación se describe cada ventana apoyándose en una captura representativa.

\subsection{Ventana \textit{Pipeline}: lanzamiento y seguimiento}
\label{anexo:ui-pipeline}

La ventana \textit{Pipeline} (Figura~\ref{fig:ui-pipeline}) es el punto de entrada
del sistema. La columna izquierda, \textit{Start Pipeline Run}, concentra los dos
únicos artefactos que el usuario debe aportar antes de ejecutar:

\begin{enumerate}
  \item un \textbf{esquema} en formato JSON (\textit{Upload schema JSON}), de
        provisión obligatoria, que declara las columnas esperadas y sus tipos y que
        gobierna la validación determinista del \textit{dataset};
  \item uno o varios \textbf{ficheros CSV} de datos (\textit{Upload CSV files}).
\end{enumerate}

El botón \textit{Run Pipeline} permanece deshabilitado hasta que se aporta el esquema,
y junto a él una etiqueta de estado (\texttt{idle} en la captura) refleja la situación
de la ejecución. La columna derecha, \textit{Events}, materializa el seguimiento en
tiempo real mediante un flujo de eventos servidos por \textit{Server-Sent Events}.
Está organizada en tres pestañas que separan tres niveles de detalle sobre una misma
traza, evitando que el registro técnico ahogue la narración de la ejecución:

\begin{itemize}
  \item \textbf{Timeline} --- hitos depurados y deduplicados de la ejecución
        (\textit{dataset} cargado, validación superada, puerta de aprobación
        solicitada, entrenamiento completado, etc.).
  \item \textbf{Tool Details} --- uso agregado de herramientas por agente y actividad
        de los nodos LLM de un solo disparo.
  \item \textbf{Raw Logs} --- la traza íntegra evento a evento, descargable como JSON
        para auditoría.
\end{itemize}

Una vez en marcha, esta misma ventana incorpora una cabecera de ejecución y un
\textit{stepper} horizontal de ocho etapas que distingue visualmente el tipo de cada
nodo (agente, LLM, determinista o puerta humana), de manera que el usuario identifica
la fase activa sin necesidad de leer los registros. Es también aquí donde se
renderizan las dos puertas de aprobación humana (Sección~\ref{sec:hitl-gates}): la
revisión del \textit{dataset} y la autorización del despliegue.

\begin{figure}[H]
\centering
% Fuente: captura de la ventana Pipeline del panel de control
\includegraphics[width=\textwidth]{figuras/anexo_ui/pipeline.png}
\caption{Ventana \textit{Pipeline} en estado inicial (\texttt{idle}). A la izquierda,
 el formulario de lanzamiento exige un esquema JSON antes de habilitar la ejecución y
 admite la carga de ficheros CSV; a la derecha, el panel \textit{Events} ---con sus
 pestañas \textit{Timeline}, \textit{Tool Details} y \textit{Raw Logs}--- queda a la
 espera de eventos.}
\label{fig:ui-pipeline}
\end{figure}

\subsection{Ventana \textit{Experiments}: consulta de ejecuciones}
\label{anexo:ui-experiments}

La ventana \textit{Experiments} (Figura~\ref{fig:ui-experiments}) ofrece una vista de
una ejecución cada vez sobre el historial registrado en MLflow. Un desplegable
superior selecciona el experimento (\texttt{mlops-agents} en la captura) y la barra
lateral lista las ejecuciones que contiene ---distintas familias de modelo de
pronóstico: \texttt{lightgbm\_forecaster}, \texttt{extra\_trees\_forecaster},
\texttt{random\_forest\_forecaster}, \texttt{xgboost\_forecaster},
\texttt{auto\_arima}, etc.--- resaltando con el color de acento la seleccionada.

El panel principal es un \emph{detalle de ejecución sin gráficas}: una decisión de
diseño deliberada. Las gráficas previas (curvas de entrenamiento, radar, barras)
visualizaban siempre el mismo uno o dos valores de una única ejecución ---los modelos
de pronóstico registran cada métrica una sola vez--- por lo que se sustituyeron por
una lectura directa y honesta de los datos que la API ya devuelve. El panel se
estructura en tres partes:

\begin{enumerate}
  \item \textbf{Cabecera} --- nombre de la ejecución (\texttt{extra\_trees\_forecaster}),
        una insignia de estado (\textit{complete}), la marca temporal, el
        \texttt{run\_id} truncado en tipografía monoespaciada y el botón
        \textit{Export CSV}, que exporta métricas y parámetros en formato RFC~4180.
  \item \textbf{Metrics} --- una rejilla de tarjetas, una por métrica, ordenadas
        alfabéticamente (en la captura, \texttt{rmse}~$=91{,}563$).
  \item \textbf{Configuration} --- los hiperparámetros de la ejecución como tabla
        clave/valor (\texttt{lags}~$=23$, \texttt{max\_depth}~$=5$,
        \texttt{n\_estimators}~$=89$).
\end{enumerate}

Esta ventana es, en la práctica, la lente sobre el seguimiento de experimentos: cada
candidato entrenado por el ejecutor (Sección~\ref{sec:nodo-executor}) queda aquí
consultable con sus métricas y su configuración exacta.

\begin{figure}[H]
\centering
% Fuente: captura de la ventana Experiments del panel de control
\includegraphics[width=\textwidth]{figuras/anexo_ui/experiments.png}
\caption{Ventana \textit{Experiments} con la ejecución \texttt{extra\_trees\_forecaster}
 seleccionada. La barra lateral enumera las familias de modelo registradas en el
 experimento \texttt{mlops-agents}; el panel de detalle muestra la métrica
 (\texttt{rmse}) y la configuración (\texttt{lags}, \texttt{max\_depth},
 \texttt{n\_estimators}), con exportación a CSV.}
\label{fig:ui-experiments}
\end{figure}

\subsection{Ventana \textit{Observability}: salud, coste y uso de herramientas}
\label{anexo:ui-observability}

La ventana \textit{Observability} (Figura~\ref{fig:ui-observability}) agrega la
telemetría operativa del sistema en cuatro tarjetas. Un principio rige su diseño:
\emph{nunca mostrar datos que no se poseen}; los apartados sin información presentan un
estado vacío explícito en lugar de inventar cifras.

\begin{itemize}
  \item \textbf{Pipeline health} --- recuento de las últimas ejecuciones desde el
        arranque del servidor, clasificadas en exitosas, fallidas y a la espera de
        intervención humana. La tarjeta advierte honestamente que el contador es en
        memoria y \textit{se reinicia con el contenedor}, pues no se persiste el
        historial en esta iteración.
  \item \textbf{LLM activity (current run)} --- actividad de los nodos LLM de la
        ejecución en curso (llamadas, duración y estado). En la captura aparece vacía
        porque no hay ejecución activa.
  \item \textbf{LLM Usage \& Estimated Cost} --- consumo de \textit{tokens} y coste
        estimado a partir de una tabla de precios local (\texttt{model\_pricing.yaml}):
        \texttt{gpt-5.4-mini} a \$0{,}75/\$4{,}50 y \texttt{gpt-5.4-nano} a
        \$0{,}20/\$1{,}25 por millón de \textit{tokens} de entrada/salida. La nota al
        pie deja claro que el coste es una \emph{estimación} basada en precios locales
        actualizables, no una facturación real.
  \item \textbf{Tool usage (current run)} --- uso agregado de las herramientas
        invocadas por los agentes en la ejecución en curso.
\end{itemize}

Esta ventana cierra el relato del sistema al hacer tangible el coste de la autonomía
agéntica ---el desglose analizado en la Sección~\ref{sec:coste-agente}--- y el reparto
del trabajo entre razonamiento LLM y ejecución determinista.

\begin{figure}[H]
\centering
% Fuente: captura de la ventana Observability del panel de control
\includegraphics[width=\textwidth]{figuras/anexo_ui/observability.png}
\caption{Ventana \textit{Observability} sin ejecución activa. \textit{Pipeline health}
 resume las últimas ejecuciones (señalando que se reinicia al reiniciar el contenedor);
 las tarjetas de actividad y coste LLM, y de uso de herramientas, muestran estados
 vacíos explícitos en lugar de cifras inventadas. La estimación de coste se apoya en
 una tabla de precios local.}
\label{fig:ui-observability}
\end{figure}
